\section{Eigenvector Continuation}
Eigenvector continuation is a reduced basis technique that allows us to find the eigenvectors and eigenvalues of a Hamiltonian that depends on a parameter $\lambda$ at a target $\lambda_*$ by projecting the Hamiltonian at $\lambda_*$ onto a subspace spanned by eigenvectors of the Hamiltonian at neighboring values of $\lambda$. This can reduce the size of the effective Hamiltonian significantly, making the problem more computationally tractable.
\subsection{Motivation of Eigenvector Continuation}
\begin{itemize}
	\item Suppose we have a parameter-dependent Hamiltonian:
	\[
	H(c) = H_0 + c H_1
	\]
	\item In practice, direct diagonalization of $H(c^*)$ at the target value $c^*$ may be too expensive (e.g. large lattice).
\end{itemize}

\subsection{Eigenvector Continuation (EC): Core Idea}
\begin{itemize}
	\setlength{\itemsep}{5pt}
	\item Explicitly diagonalize $H(c_i)$ at several tractable sample points $\{c_i\}$.
	\item Collect the corresponding eigenvectors $\{ |v_i\rangle \}$.
	\item Approximate the target state $|v(c^*)\rangle$ as a linear combination:
	\[
	|v(c^*)\rangle \approx \sum_i c_i |v_i\rangle
	\]
	\item Project $H(c^*)$ onto the subspace spanned by $\{ |v_i\rangle \}$:
	\[
	\sum_j \langle v_i|H(c^*)|v_j\rangle c_j 
	= E \sum_j \langle v_i|v_j\rangle c_j
	\]
	\item Solve the generalized eigenvalue problem:
	\[
	H_{\text{proj}} c = E S c
	\]
\end{itemize}

\begin{itemize}
	\item This only assumes that the eigenvectors at the target parameter $c$ lie in the span of the sampled eigenvectors.
	\item If $H(c)$ is analytic in $c$, its eigenvalues and eigenvectors are also analytic (from perturbation theory), so locally: $|\psi(c)\rangle = |\psi(c_0)\rangle + \left.\frac{d|\psi\rangle}{dc}\right|_{c_0} (c - c_0) + \dots$
	\item Approximating the derivative via finite differences between neighboring points:
	\[
	\frac{d|\psi\rangle}{dc}\Big|_{c_i} \approx \frac{|\psi(c_{i+1})\rangle - |\psi(c_i)\rangle}{c_{i+1}-c_i}
	\]
	\item Interpretation: $|\psi(c)\rangle$ is approximately a linear combination of eigenvectors at neighboring sampled points.
	\item Hence, as $c$ varies, low-lying eigenvectors trace out only a small subspace of Hilbert space.
	with coefficients $\alpha_i$ determined by interpolation between neighboring points.
	\item Interpretation: the low-lying eigenvectors explore only a small portion of Hilbert space as $c$ varies; they trace out a low-dimensional subspace.
\end{itemize}

\subsubsection{Computational Gain}
\begin{itemize}
	\item If the full Hamiltonian acts on a space of dimension $10^6$ (e.g., $100^3$ lattice points),
	EC reduces the problem to diagonalizing a much smaller matrix—say $100 \times 100$.
	\item By solving a few smaller problems (for several $c_i$), we can predict eigenvalues and eigenvectors across the full parameter range of $H(c)$.
	\item EC thus provides a non-perturbative interpolation scheme for eigenstates of large Hamiltonians.
\end{itemize}


\subsection{Finite Volume Eigenvector Continuation (FVEC)}

\subsubsection{Extending EC to Finite Volume}

\begin{itemize}
	\item Consider simulations in finite periodic boxes of volumes $\{V_i\}$.
	\item Eigenvalues and eigenvectors vary smoothly with volume, encoding information about the infinite-volume limit.
	\item FVEC generalizes EC by treating the parameter as the box volume $V$ instead of $c$.
\end{itemize}


\subsubsection{Technical Subtleties of FVEC}

\begin{itemize}
	\item The inner product depends on $V$:
	\[
	\langle v_i | v_j \rangle_V = \int_V d^3x\, v_i^*(x) v_j(x)
	\]
	so we cannot directly mix vectors from different volumes.
	\item Define an enlarged Hilbert space:
	\[
	\mathcal{H} = \bigoplus_i \mathcal{H}_{V_i}
	\]
	\item States are dilated to a common reference volume $V_*$, allowing projections like
	\[
	\langle v_i | H(V^*) | v_j \rangle_*,
	\]
	where the subscript $*$ denotes the chosen common inner product.
	\item The EC algorithm then proceeds identically.
\end{itemize}



